% Requires amsmath/amssymb and theorem environments from the host manuscript.
% Add the separate bibliography entry hsd:continuation.
\section{Higher spectral defects and beta polarization}
\label{hsd:int:section}
Fix $a>0$, $W(x)=\sum_{d=0}^D w_dx^d$, and define the continued linear spectrum
\[
 F_{a,W}(s;z)=\sum_{d=0}^D w_d\sum_{h=0}^d
     \binom dh(-z)^{d-h}\zeta(s-h,a+z),\qquad \Re(a+z)>0.
\]
Let $Z_W(s;a,b)=\sum_{n\ge0}W(n+a)((n+a)^2+b^2)^{-s}$ initially in its
convergence half-plane, and set
\[
 \Delta_r(W;a,b)=2^{r-1}\{F_{a,W}^{(r)}(0;ib)+F_{a,W}^{(r)}(0;-ib)\}
                                      -Z_W^{(r)}(0;a,b).
\]
The superscript denotes spectral differentiation.
For an analytic endpoint function write
\[
 \mathcal C_t f(u,v)=f(u)+f(v)-f((1-t)u+tv)-f(tu+(1-t)v).
\]
The endpoint-subtracted beta master in \cite{hsd:continuation} gives
\begin{equation}\label{hsd:int:beta}
 \frac{F_{a,W}(2s;ib)+F_{a,W}(2s;-ib)}2-Z_W(s;a,b)
 =\frac{s\Gamma(1+2s)}{4\Gamma(1+s)^2}
 \int_0^1[t(1-t)]^{s-1}\mathcal C_tF_{a,W}(2s;\cdot)(ib,-ib)\,dt.
\end{equation}
The integral is ordinary and holomorphic near $s=0$ after the displayed
subtraction. Its derivatives may be taken under the integral.

Write $\Gamma(1+2s)/\Gamma(1+s)^2=\sum_{m\ge0}\beta_ms^m$ and
\[
 I_{j,p}=\int_0^1\frac{\log^p(t(1-t))}{t(1-t)}
          \mathcal C_t(F_{a,W}^{(j)}(0;\cdot))(ib,-ib)\,dt.
\]
Then, at every $r\ge1$,
\begin{equation}\label{hsd:int:alljets}
 \Delta_r=\frac{r!}{4}\sum_{m+j+p=r-1}\frac{2^j\beta_m}{j!p!}I_{j,p}.
\end{equation}
Thus the inner spectral order is at most $r-1$.
The zero jet is
$F_{a,W}(0;z)=\sum_dw_d\zeta(-d,a)+\sum_dw_d(-1)^{d+1}z^{d+1}/(d+1)$,
so its endpoint subtraction is a local polynomial independent of $a$.
At conjugate endpoints it vanishes for even $W$; the positive-order
integrals need not vanish.

\subsection{The exact second defect}
Let $\Phi_{a,W}(z)=F_{a,W}'(0;z)$ and
\[
 K_{a,W}(u,v)=\int_0^1\frac{\mathcal C_t\Phi_{a,W}(u,v)}{t(1-t)}\,dt,
 \qquad O_k^{(p)}=H_{2k}^{(p)}-2^{-p}H_k^{(p)},\quad O_k=O_k^{(1)}.
\]
\begin{theorem}\label{hsd:int:second}
For every polynomial $W$ and $a,b>0$,
\[
 \Delta_2(W;a,b)=K_{a,W}(ib,-ib)
 -\sum_{k\ge1}\frac{2(-1)^kw_{2k-1}b^{2k}}k
                          \bigl(O_k^2+O_k^{(2)}\bigr).
\]
The sum is finite. For fixed $a>0$, the identity
$\Delta_2(W;a,b)=0$ for every $b>0$ holds if and only if $W=0$.
\end{theorem}
The local formula follows from the second beta row and the moment
\[
 \int_0^1\frac{1-(1-2t)^{2k}}{t(1-t)}\log(t(1-t))\,dt
                     =-4\bigl(O_k^2+O_k^{(2)}\bigr).
\]
For universal vanishing, sufficiently high Taylor coefficients in $b$
are nonzero multiples of $\sum_{n\ge0}W(n+a)/(n+a)^{2k}$. The first
nonzero lattice value of a nonzero polynomial dominates these moments
as $k\to\infty$; see the full proof in \cite{hsd:continuation}.
In particular,
\[
 \Delta_2(1;a,b)=-4\sum_{n\ge0}\arctan^2\!\frac b{n+a},
\]
so even parity is not a second-order vanishing criterion.

\subsection{All quadratic masses and their coincident limit}
Assume $W(x)=P(x^2)$ and put
$T_{a,W}(\mu)=F_{a,W}'(0;i\sqrt\mu)+F_{a,W}'(0;-i\sqrt\mu)$.
Define the ordinary mass chord
\[
 V_{a,W}(b,c)=\int_0^1\frac{\mathcal C_tT_{a,W}(b^2,c^2)}{t(1-t)}\,dt.
\]
For $Z_{W;\mathbf b}(s;a)=\sum_{n\ge0}W(n+a)
\prod_{j=1}^r((n+a)^2+b_j^2)^{-s}$, continuation gives
\begin{equation}\label{hsd:int:masses}
 Z_{W;\mathbf b}''(0;a)=r\sum_{j=1}^r Z_W''(0;a,b_j)
                         -\sum_{i<j}V_{a,W}(b_i,b_j).
\end{equation}
For $|h|<a^2+b^2$ in the analytic squared-mass continuation,
\begin{equation}\label{hsd:int:confluence}
 V_{a,W}(b,\sqrt{b^2+h})=
 \sum_{k\ge2}\frac{2(-1)^kH_{k-1}}k h^kZ_W(k;a,b).
\end{equation}
All positive-integer coefficients reduce finitely to polygamma functions
and Bernoulli polynomials. For example
$Z_1(1;a,b)=(\psi(a+ib)-\psi(a-ib))/(2ib)$; higher values follow by
mass differentiation and polynomial division of $W$.
The diagonal is regular and $V(b,b)=0$.

On the bilateral lattice, replace $T_{a,W}$ by
\[
 T_{a,W}^{\rm bil}(\mu)=\sum_{\ell=0}^{\deg P}
              \frac{P^{(\ell)}(-\mu)}{\ell!}L_\ell(a,\sqrt\mu),
 \qquad L_\ell(a,b)=E'(-\ell;a,b).
\]
The established finite polylogarithm formulas for $L_\ell$ thus give
finite polylogarithmic covariance integrands. The full report also gives
an explicit finite-integrand identity for every $E''(-m;a,b)$.

\subsection{Normalized mass primitives}
For $E(s;a,b)=\sum_{n\in\mathbb Z}((n+a)^2+b^2)^{-s}$, $0<a<1$, define
\[
 A_{m,r}(a,\mu)=\left.\partial_t^r\left\{
                \frac{m!}{(1-t)_m}E(t-m;a,\sqrt\mu)\right\}\right|_{t=0}.
\]
Then $\partial_\mu A_{m,r}=mA_{m-1,r}$ and
\[
 A_{m,r}(a,\mu)=\sum_{j=0}^{m-1}\binom mj\mu^j A_{m-j,r}(a,0)
       +m\int_0^\mu(\mu-u)^{m-1}E^{(r)}(0;a,\sqrt u)\,du.
\]
All boundary values are finite Hurwitz jets. At second order
$A_{m,2}=E''(-m)+2H_mE'(-m)$; omitting this harmonic normalization gives
an inhomogeneous, not a homogeneous, ladder.

These results settle explicit representation and confluence questions.
They do not establish a smallest function algebra or irreducibility of
the remaining ordinary integrals. Classical attribution and the
canonical S6/revised-S8 statuses remain unchanged.
